\chapter{Preferential Two-Stage Selection from Open Law Classes}
\label{ch:open-formation}

This chapter concerns the reciprocal law potential of \cref{def:dyn-potential} with exact Gibbs rows (\cref{prop:dyn-gibbs}); no statement below is transferred to the per-agent or global free energies of \cref{ch:dynamics}. The common-precision construction selects an exact additive description within a prepared family. The construction given here instead admits initial laws ranging over a full open neighborhood. Native source rows, read by a fixed threshold selector, select three pairs and then a pair of their parents. The selection reads only a transmitted bit. The compression is an orientation symmetry of a fixed selective channel; it exists before the selected events and descends for every partition at every time. What the events change is the resolution of a declared external response interface, and the information discarded at a handoff is exactly the information that this interface stops resolving there. Measured against the formation standard \cref{inh:recursive-formation}, the main theorem establishes selection, exact inheritance, and compression on an open class over a finite window. It does not establish interaction necessity (F5): in the regime of its proof, integrations indicate that the null flow of \cref{inh:null-flow} produces the same membership history (\cref{open:null-control}). Every block relation is transient, and the selector later merges the exterior pair (\cref{open:final-merge}). The construction is therefore a two-stage selection with exact inheritance, not a recursive formation.

\section{A homogeneous selective-channel model}

Every primitive belief and model space is $\mathsf K=\mathsf M=\{-1,1\}\times S^1$. Use the probability measure
\begin{equation}
d\nu=\frac12\sum_{z=\pm1}\delta_z\otimes\frac{d\theta}{2\pi}.
\end{equation}
All densities below are relative to this measure. The $2N$ normalized densities $q_i,s_i$ vary independently over positive densities of class $C^r$ in the angle, with $r\ge3$ and logarithms bounded in $C^r$; this is the chart of the chapter. Internal bit-angle correlations are unrestricted. Recognition is the declared product $q_i\otimes s_i$. Take uniform proper references and evaluators, $r_i=G_i=1$, equal unit Fisher mobilities, a fixed temperature $T>0$, and all-source priors $\pi_{ij}=1/N$, including self.

Use the same fixed normalized comparison channel in both sectors and on every edge, self edges included:
\begin{equation}
\mathcal K((z,\theta),d(z',\theta'))
=\delta_z(dz')\frac{d\theta'}{2\pi}.
\label{open:channel}
\end{equation}
It transmits the bit and replaces the angle by an independent uniform draw. There is no channel switch at a selected event. In the terminology of \cref{ch:smooth-closure}, the bit is the public coordinate, meaning the coordinate that peer channels transmit, and the angle is a private coordinate with Haar private reference. The word public is used only in this sense below. The channel is the $S^1$ twirl of the angular coordinate, which averages a law over all rotations of its angle; it is the map that describes communication between parties that share no angular reference frame \citep{BartlettRudolphSpekkens2007}.

If
\begin{equation}
\mu_i(z)=\int q_i(z,\theta)\frac{d\theta}{2\pi},
\qquad
\eta_i(w)=\int s_i(w,\psi)\frac{d\psi}{2\pi},
\end{equation}
then $q_j\mathcal K$ has density $\mu_j(z)$, and the model output has density $\eta_j(w)$. The chain rule gives
\begin{equation}
\KL(q_i\|\mu_j\otimes\mathrm{Haar})
=H_i^q+\KL(\mu_i\|\mu_j),
\qquad
H_i^q=\KL(q_i\|\mu_i\otimes\mathrm{Haar}).
\label{open:cost}
\end{equation}
The receiver term $H_i^q$ cancels from the Gibbs normalization. Thus
\begin{equation}
\beta_{ij}
=\frac{\exp[-\KL(\mu_i\|\mu_j)/T]}
{\sum_k\exp[-\KL(\mu_i\|\mu_k)/T]},
\label{open:row}
\end{equation}
and $\gamma$ is the same expression in $\eta$. All entries are positive. The cancellation uses the selective self channel. With the identity self channel of \cref{def:dyn-potential} one would have $D^q_{ii}=0$ while every peer cost retains $H_i^q$, so each peer weight would carry the factor $e^{-H_i^q/T}$ relative to self; at the small temperatures used below this suppresses every peer weight unless $H_i^q$ is of order $T$. The construction therefore requires the selective self channel, which is a constitutive choice. The angular term remains in the reduced energy:
\begin{equation}
\mathcal E_{\rm red}
=\sum_i\bigl[\KL(q_i\|1)+\KL(s_i\|1)\bigr]
-T\sum_i\log\left(\frac1N\sum_j e^{-D^q_{ij}/T}\right)
-T\sum_i\log\left(\frac1N\sum_j e^{-D^s_{ij}/T}\right),
\end{equation}
where $D^q_{ij}$ is the left side of \cref{open:cost}.

\propositionheading{Complete density force}{open:force}
On the positive chart, the exact unit-Fisher field is
\begin{equation}
\begin{split}
h_i^q(z,\theta)
&=2\log q_i(z,\theta)-\sum_j\beta_{ij}\log\mu_j(z)
-\sum_k\beta_{ki}\frac{\mu_k(z)}{\mu_i(z)},\\
\partial_tq_i&=-q_i\bigl(h_i^q-\E_{q_i}h_i^q\bigr),
\end{split}
\label{open:full-field}
\end{equation}
with the exchanged formula for $s,\eta,\gamma$. Every sum includes self.
\begin{proof}
The local KL contributes $\log q_i$ modulo a constant. Differentiating a reduced row by its envelope identity gives its weighted divergence derivatives. The normalized receiver row contributes another $\log q_i-\sum_j\beta_{ij}\log\mu_j$. Varying a source marginal in every incoming row gives $-\sum_k\beta_{ki}\mu_k/\mu_i$. The fixed channel has no law derivative. Fisher inversion on normalized densities subtracts the $q_i$ expectation and multiplies by $-q_i$.

For a self term the incoming ratio is one, but $D^q_{ii}=H_i^q$ need not vanish. Its constant incoming contribution and its angular entropy are both included. Logarithms and marginal ratios are smooth on a neighborhood bounded away from zero. Bounded integration, multiplication, and finite Gibbs normalization preserve $C^r$. The field contains no angular derivative, so it is locally Lipschitz both for the $C^r$ norm of log-densities and for their supremum norm, and it has a unique local interior flow on the chart. Its centering preserves normalization.
\end{proof}

\section{Two stages selected by native routing}

For a current partition $\mathcal P$, retain the actual block rows
\begin{equation}
B^q_{IJ}=\frac1{|I|}\sum_{i\in I}\sum_{j\in J}\beta_{ij},
\qquad
B^s_{IJ}=\frac1{|I|}\sum_{i\in I}\sum_{j\in J}\gamma_{ij}.
\label{open:block-row}
\end{equation}
Self mass remains in each normalized row. Its experiment chooses a receiver uniformly in $I$, draws its native source index, and reports the source block. If blocks are joined, finite summation gives
\begin{equation}
B^a_{CD}
=\frac1{|C|}\sum_{I\subset C}|I|\sum_{J\subset D}B^a_{IJ},
\qquad a\in\{q,s\}.
\label{open:row-recursion}
\end{equation}
This is associative aggregation with inherited primitive counts.

\definitionheading{Preferential pair rule}{open:selector}
Fix $\varepsilon=1/10$ and $\delta=1/20$. Distinct current blocks $I,J$ qualify when, in both sectors and directions,
\begin{equation}
B^a_{IJ}>\varepsilon,
\qquad
B^a_{IJ}-\sum_{K\in\mathcal P\setminus\{I,J\}}B^a_{IK}>\delta.
\label{open:contrast}
\end{equation}
The exterior excludes self and the proposed partner. A partner need not beat the self loop. Merge qualifying disjoint pairs at their first closed-threshold contact when the strict inequalities hold immediately afterward: the event time is the first time at which the closed inequalities hold and the strict inequalities hold on a right neighborhood. In every event proved below the relevant crossing is transverse, so this time is the unique zero of a block-row difference with positive time derivative. A block cannot have two qualifying partners: each would have to exceed the other by more than $\delta$. The rule is therefore unambiguous at simultaneous disjoint events and respects relabeling. Persistence below means strict contrast over the stated positive time windows.

The routing calculation first uses binary densities $\mu_i(z)=1+zx_i$, $\eta_i(w)=1+wy_i$, with $(x,y)\in(-1,1)^{2N}$. Define $d(u,v)=\KL(\mu_u\|\mu_v)$. Its derivatives are
\begin{equation}
\partial_ud(u,v)=\operatorname{atanh}u-\operatorname{atanh}v,
\qquad
\partial_vd(u,v)=\frac{v-u}{1-v^2}.
\end{equation}
The exact binary field is
\begin{equation}
\begin{split}
\dot x_i={}&(1-x_i^2)\left[-\operatorname{atanh}x_i
-\sum_j\beta_{ij}(\operatorname{atanh}x_i-\operatorname{atanh}x_j)\right]\\
&+\sum_k\beta_{ki}(x_k-x_i).
\end{split}
\label{open:binary-field}
\end{equation}
The model equation is identical and independent. Self summands vanish in this binary equation, consistently with $d(x_i,x_i)=0$.

\lemmaheading{Interior binary evolution}{open:interior}
Every binary initial configuration has a unique global interior solution, and all its means tend to zero.
\begin{proof}
At a maximum of $\vartheta_i=\operatorname{atanh}x_i$, both peer sums in the equation for $\dot\vartheta_i$ are nonpositive; at a minimum they are nonnegative. The finite maximum and minimum satisfy the corresponding one-sided differential inequalities. Scalar comparison yields
\begin{equation}
e^{-t}\min_i\vartheta_i(0)\le\vartheta_i(t)
\le e^{-t}\max_i\vartheta_i(0).
\label{open:logit-bound}
\end{equation}
The same holds for the model logits. The solution stays in a compact interior box, where the smooth field extends it past any proposed finite maximal time. The bounds imply convergence. This does not by itself establish the order of selection events.
\end{proof}

The perturbation steps below all use one elementary statement. It compares a reference trajectory with a perturbed one and tracks both the trajectory and the observables that define events, since the observables themselves change with each perturbation.

\lemmaheading{Persistence of strict windows}{open:persistence}
Let $X$ solve $\dot X=\Phi(X)$ on $[0,L]$ in a finite-dimensional space with the supremum norm, and let $K_\rho$ be the tube of closed balls of radius $\rho$ about $X(t)$. On $K_\rho$ suppose that a possibly time-dependent field $\Psi$ satisfies $|\Psi-\Phi|\le e$, that $\Phi$ is $C$-Lipschitz on each slice, and that $|\Phi|,|\Psi|\le M$. Let $O$ be a $C^1$ observable with $|\nabla O|_1\le G$ and $\nabla O$ Lipschitz with constant $L_1$, and let $\widetilde O$ satisfy $|\widetilde O-O|\le\eta_0$ and $|\nabla\widetilde O-\nabla O|_1\le\eta_1$ on $K_\rho$. If $\dot Y=\Psi(Y)$, $|Y(0)-X(0)|\le\delta_0$, and
\begin{equation}
\Delta=\delta_0e^{CL}+e\frac{e^{CL}-1}{C}<\rho,
\end{equation}
then $Y(t)\in K_\rho$ on $[0,L]$, and
\begin{equation}
\begin{aligned}
|\widetilde O(Y)-O(X)|&\le\eta_0+G\Delta,\\
\Bigl|\tfrac{d}{dt}\widetilde O(Y)-\tfrac{d}{dt}O(X)\Bigr|
&\le\eta_1M+L_1\Delta M+G(e+C\Delta).
\end{aligned}
\label{open:persistence-bounds}
\end{equation}
Suppose $|O(X)-\theta_0|\ge m$ outside an interval $I$, $O(X)-\theta_0$ has opposite signs at the ends of $I$, and $\frac{d}{dt}O(X)\ge s>0$ on $I$. If the two bounds in \cref{open:persistence-bounds} are below $m$ and $s$ respectively, then $\widetilde O(Y)-\theta_0$ has exactly one zero in $[0,L]$, it lies in $I$, and it is transverse; strict windows of the reference persist.
\begin{proof}
While $Y$ is in the tube, $|Y-X|'\le|\Psi(Y)-\Phi(Y)|+|\Phi(Y)-\Phi(X)|\le e+C|Y-X|$ because each slice is convex. Gronwall's inequality bounds $|Y-X|$ by $\Delta<\rho$, so the first exit time cannot occur. The first observable bound follows from $|\widetilde O(Y)-O(Y)|+|O(Y)-O(X)|$ and the mean value theorem on the convex slice. For the second, write $\frac{d}{dt}\widetilde O(Y)=\nabla\widetilde O(Y)\cdot\Psi(Y)$ and $\frac{d}{dt}O(X)=\nabla O(X)\cdot\Phi(X)$, and split the difference into the four terms $(\nabla\widetilde O-\nabla O)(Y)\cdot\Psi(Y)$, $(\nabla O(Y)-\nabla O(X))\cdot\Psi(Y)$, $\nabla O(X)\cdot(\Psi-\Phi)(Y)$, and $\nabla O(X)\cdot(\Phi(Y)-\Phi(X))$. Outside $I$ the perturbed observable keeps the reference sign; inside $I$ it is strictly increasing and changes sign once.
\end{proof}

\propositionheading{Two-stage preferential routing}{open:routing}
For $N=6$, some fixed $T>0$ and a nonempty relabeling-invariant open set in $(-1,1)^{12}$ have the following behavior. Every initial off-diagonal native weight is below $1/10$. Three preferential pairs form and coexist for a positive interval. Two of their parents subsequently join while the third remains an exterior with strictly smaller routing masses. The first pairs and the selected parent pair satisfy \cref{open:contrast} over positive windows. All events occur at finite interior times under \cref{open:binary-field}, with no parameter change.
\begin{proof}
Choose $T=h^2$ and write $x_i=hu_i$. The function $h^{-2}d(hu,hv)$ is even and analytic in $h$ on $h\max(|u|,|v|)<1$, with
\begin{equation}
h^{-2}d(hu,hv)=\frac{(u-v)^2}{2}+\frac{h^2}{12}(u-v)^2(u^2+2uv+3v^2)+O(h^4).
\end{equation}
Dividing \cref{open:binary-field} by $h$, without changing time, gives a field and block rows that are analytic in $(h^2,u)$ on $\{h|u|_\infty<1\}$. At $h=0$ the field is
\begin{equation}
\dot u_i=-u_i+\sum_j(w_{ij}+w_{ji})(u_j-u_i),
\qquad
w_{ij}=\frac{e^{-(u_i-u_j)^2/2}}{\sum_k e^{-(u_i-u_k)^2/2}}.
\label{open:limit}
\end{equation}
The self summand one bounds every denominator away from zero. Limit trajectories remain bounded on finite intervals by the same extrema comparison. The limit \cref{open:limit} is the Euclidean gradient flow of the Gaussian log-sum-exp energy $\sum_iu_i^2/2-\sum_i\log\sum_ke^{-(u_i-u_k)^2/2}$. Its pair formation is the clustering mechanism of Gaussian mean shift and its blurring variant \citep{Cheng1995,CarreiraPerpinan2006} and of self-attention particle dynamics \citep{GeshkovskiEtAl2025}; here it is obtained from the original KL source rows at low temperature. The rest of the proof constructs a reference trajectory of \cref{open:limit} with strict windows and transfers them three times by \cref{open:persistence}: from decoupled clusters to the six-agent limit field, from the limit field to the exact scaled field, and from symmetric to open initial data.

For four agents first take the reflection-invariant configuration
\begin{equation}
u=(-a-b,-a+b,a-b,a+b),\qquad a(0)=A,\quad b(0)=2.
\end{equation}
Here $a$ is half the distance between the two pair centers and $b$ is the internal half-gap. Set $w=e^{-2b^2}$, $v_0=e^{-2a^2}$, $v_\pm=e^{-2(a\pm b)^2}$,
$Z_\pm=1+w+v_0+v_\pm$, and
\begin{equation}
c_p=w(Z_+^{-1}+Z_-^{-1}),
\qquad c_0=v_0(Z_+^{-1}+Z_-^{-1}).
\end{equation}
Substitution into the full field gives
\begin{equation}
\begin{split}
\dot a&=-a-2ac_0
-2\left[\frac{(a+b)v_+}{Z_+}+\frac{(a-b)v_-}{Z_-}\right],\\
\dot b&=-b-2bc_p
-2\left[\frac{(a+b)v_+}{Z_+}-\frac{(a-b)v_-}{Z_-}\right].
\end{split}
\label{open:core}
\end{equation}
Both unequal row normalizers are retained.

For $a\ge4$ and $0\le b\le2$, put $c=1+w+v_0\ge1$. The function
$g_c(r)=r e^{-2r^2}/(c+e^{-2r^2})$ satisfies
\begin{equation}
|g_c'(r)|\le(1+4r^2)e^{-2r^2}.
\end{equation}
This bound decreases for $r\ge2$, so
$|g_c(a+b)-g_c(a-b)|\le34be^{-8}$.
Since $0\le c_p\le1$, $2ac_0\le4ae^{-32}$, and $(a+b)v_++(a-b)v_-\le2av_-\le2ae^{-8}$, \cref{open:core} implies
\begin{equation}
-(3+68e^{-8})b\le\dot b\le-(1-68e^{-8})b,\qquad -(1+4e^{-32}+4e^{-8})a\le\dot a\le-a.
\label{open:early-bounds}
\end{equation}
Here $1-68e^{-8}>0.977$ and $1+4e^{-32}+4e^{-8}<1.0014$. The positive coordinates cannot first reach their invariant zero axes at finite time.

Take $A\ge34$. The first time $\tau$ with $b=1/2$ satisfies
$\tau\le\log4/0.977<1.42$, and $a(t)\ge Ae^{-1.0014\cdot1.42}>0.239A\ge8$ through this interval. Initial peer weights are at most $e^{-8}$ and cross weights much smaller. A directed within-pair weight is $p_\pm=w/Z_\pm$. While $b\ge1/2$,
$\dot w\ge w/2$, whereas
\begin{equation}
|\dot v_0+\dot v_\pm|
\le8a^2e^{-2a^2}
+4(a+2)(2a+8)e^{-2(a-2)^2}<1/100
\end{equation}
for $a\ge8$. Both positive terms decrease on that range. Differentiating $p_\pm$ therefore gives $\dot p_\pm>0$. At $\tau$,
\begin{equation}
p_\pm\ge\frac{e^{-1/2}}{1+e^{-1/2}+2e^{-8}}>1/4.
\end{equation}
Each peer threshold is crossed transversely, with exterior mass less than $2e^{-8}$. The contrast is strict at entry.

Until the first time $t_A$ with $a=4$, the pair weights remain above $1/4$ and every cross weight is at most $e^{-8}$. An unwanted parent source column has at most two such entries, excluding a premature join even if one pair has formed before the other. Moreover,
\begin{equation}
\frac{\log(A/4)}{1.0014}\le t_A\le\log(A/4),
\qquad b(t_A)\le2(A/4)^{-0.975}.
\label{open:small-gap}
\end{equation}
Thus the pairs coexist while their residual internal gaps become arbitrarily small.

Starting at $a=4$, compare with the invariant trajectory $b=0$. Its exact inherited cross mass and primitive peer weight are
\begin{equation}
B=\frac{e^{-2a^2}}{1+e^{-2a^2}},\qquad
p=\frac{1-B}{2},\qquad
\dot a=-a(1+4B).
\end{equation}
Here $\dot B=4a^2e^{-2a^2}(1+4B)/(1+e^{-2a^2})^2>0$. The parent threshold is crossed at $a=\sqrt{\log9/2}$. Stop this comparison when $B=1/5$, at a finite time $S\le\log(4/\sqrt{\log4/2})$. Throughout it, $p\ge2/5$ and the primitive exterior mass is at most $1/5$, so first-pair contrast is at least $1/5$. The parent crossing has a positive derivative. By \cref{open:persistence}, with $\Psi=\Phi$ given by \cref{open:core} and the block rows as observables, these compact margins and the transverse crossing persist when the initial half-gap at $a=4$ is below some $\varepsilon_b>0$. This $\varepsilon_b$ does not depend on $A$, since the comparison starts at $a=4$. Choose one finite $A\ge34$ with $2(A/4)^{-0.975}<\varepsilon_b$ using \cref{open:small-gap}, and fix $L=t_A+S$. The parent event then occurs between $\log(A/4)/1.0014$ and $\log(A/4)+S$: the second stage happens when the supplied core separation has contracted to order one.

A third pair makes the second event have a genuine exterior. Add initial coordinates $R-2,R+2$. In the decoupled comparison the two-agent interaction is antisymmetric, so their center is $Re^{-t}$, and the half-gap $b_c$ obeys
\begin{equation}
\dot b_c=-b_c(1+4p_c),\qquad
p_c=\frac{e^{-2b_c^2}}{1+e^{-2b_c^2}},\qquad b_c(0)=2.
\end{equation}
Its threshold is transverse, and by $\log4$ its peer weight exceeds $1/4$.

Let $X(t)$ be the decoupled reference, the four-agent core trajectory together with this two-agent trajectory, and compare it in the original coordinates with the six-agent limit field. On the radius-one supremum-norm tube about $X$, the core coordinates satisfy $|u_i|\le A+3$, because the extrema comparison keeps the reference core in $[-A-2,A+2]$, and the pair coordinates lie within $3$ of $Re^{-t}$. Every core-pair separation is therefore at least
$D_R=Re^{-L}-A-6$. Let $\varphi(D)=e^{-D^2/2}$. For an agent in one cluster, the six-agent row differs from the decoupled row through the changed normalizer, $w_{ij}-\widetilde w_{ij}=-\widetilde w_{ij}\epsilon_i/Z_i$ with $\epsilon_i$ the cross-cluster numerator mass, in both the outgoing and incoming directions, and through the incoming and outgoing cross forces. Using $\widetilde w|x|\le e^{-1/2}$ and $|x|\varphi(|x|)\le D\varphi(D)$ for $|x|\ge D\ge1$, these contributions give
\begin{equation}
|\Psi-\Phi|_\infty\le8\bigl(e^{-1/2}+D_R\bigr)\varphi(D_R).
\label{open:distant-bound}
\end{equation}
The same quotient-rule bookkeeping bounds the differences of the Jacobians and of the block-row gradients by absolute multiples of $(1+D_R^2)\varphi(D_R)$; for the core-parent row, $|B^{(6)}-B^{(4)}|\le2\varphi(D_R)$. None of these constants depends on $A$ or $R$. The decoupled field is the limit field on each component. Its interaction depends only on differences, and $w|x|$ and $wx^2$ are bounded, so its Lipschitz constant on the tube is bounded independently of $A$ and $R$. Choosing $R$ sufficiently large and finite, \cref{open:persistence} preserves every window and crossing of the reference and keeps every aggregated mass between the core and the third pair below $10^{-3}$ through $L$. At the core-parent crossing the contrast is then greater than $0.099>1/20$; the third pair cannot join either core block or their union.

Fix this $R$ and let $X$ now denote the six-agent limit trajectory on $[0,L]$. On the radius-$1/2$ tube about it, $|u|_\infty\le R+3$. For $h(R+3)<1$, analyticity in $(h^2,u)$ makes the exact scaled field and its block rows converge to the limit field and its rows in $C^1$, uniformly on this compact tube, as $h\to0$. \Cref{open:persistence} with $\delta_0=0$ gives one sufficiently small fixed $h>0$ for which the exact scaled binary field has the same strict windows and transverse crossings. Choose equal sector trajectories only as a reference, then perturb all twelve initial coordinates independently. \Cref{open:persistence} with $\Psi=\Phi$ preserves the windows and the positive crossing derivatives for sufficiently small initial perturbations. Each pair's event is the latest of its four directional-sector crossings, before the next separated window. The means can be initially distinct; their invariant equality diagonals and uniqueness preclude a first finite collision. Finally take the union of all population permutations of this open neighborhood. The constitutive data contain no block labels; the initial configuration chooses the eventual blocks through its scale hierarchy.
\end{proof}

\remarkheading{Null-flow control}{open:null-control}
The null flow of \cref{inh:null-flow} deletes every peer cost. Here it keeps the angular self terms and gives, on the binary family, the decoupled field $\dot x_i=-(1-x_i^2)\operatorname{atanh}x_i$, whose scaled limit is $\dot u=-u$, so $u(t)=e^{-t}u(0)$. Read by the same rows and \cref{open:selector}, its membership history is a function of the initial positions alone: the rows at time $t$ are a Gaussian kernel of bandwidth $e^t$ on $u(0)$. Integrations of \cref{open:limit} and of this null limit from the initial data of the proof, at $A=8$, $64$, and $4096$ with a distant third pair, produce the same three pairs, the same core-parent join with the same exterior, and the same order of events; the pair events move from $t\approx0.589$ to $0.647$, and at $A=4096$ the parent join moves from $t\approx8.214$ to $8.271$. These are integrations of the limit field, not of the exact scaled field at the fixed $h$ of the proof, and they are not part of any statement proved here. What is proved is weaker and one-sided: while $a\ge4$ the internal half-gap contracts along the full field at a rate between $0.977$ and $3.03$ by \cref{open:early-bounds}, against exactly $1$ along the null limit, so the peer forces tighten the selected pairs. Neither interaction necessity (F5) nor the amplification grade of \cref{inh:recursive-formation} is established by \cref{open:routing}.

\remarkheading{Size of the constants}{open:constants}
The witness is nonconstructive. With $A=4096$ the core reference has $t_A\approx6.93$ and $L\approx8.36$, and a third pair separated by $D_R$ of order ten requires $R\approx1.8\times10^7$. Admissibility $h|u|_\infty<1$ alone then forces $T$ below about $3\times10^{-15}$. The tube step as written is more restrictive: for a far agent the scaled field is $-h^{-1}(1-h^2u^2)\operatorname{atanh}(hu)=-u+\tfrac23h^2u^3+O(h^4u^5)$, so the uniform field error on the tube contains the common-mode term $\tfrac23h^2R^3$, and keeping the far pair within a radius-$1/2$ tube needs $T=h^2$ below about $10^{-21}$ before any Gronwall amplification. Existence is unaffected, since $h$ is only required to be sufficiently small. A comparison relative to each cluster's exact center would remove the common-mode error and need only $hR$ small; we do not carry it out. The early phase above works from $A\ge34$; a numerical interval enclosure of the late comparison indicates that $\varepsilon_b=1$ suffices, which would make $A=34$ admissible, but that enclosure is not part of the proof. The temperature of the theorem below is not quantified.

\section{A full open domain of correlated laws}

The binary result embeds exactly into \cref{open:full-field}. The embedding extends to every correlated law, because the conditional angular laws evolve autonomously.

\propositionheading{Exact reduction of the bit-circle flow}{open:reduction}
Write $q_i(z,\theta)=\mu_i(z)f_i(\theta\mid z)$ and $s_i(w,\psi)=\eta_i(w)g_i(\psi\mid w)$ for positive laws with bounded log-densities. Along \cref{open:full-field},
\begin{equation}
f_i(\theta\mid z,t)=
\frac{f_i(\theta\mid z,0)^{e^{-2t}}}
{\int f_i(\psi\mid z,0)^{e^{-2t}}d\psi/(2\pi)},
\label{open:conditional-flow}
\end{equation}
for each $z$ separately, and the bit means obey \cref{open:binary-field} plus the exogenous forcing
\begin{equation}
-\frac12(1-x_i^2)\bigl[c_i(1,t)-c_i(-1,t)\bigr],
\qquad
c_i(z,t)=2\KL\bigl(f_i(\cdot\mid z,t)\big\|\mathrm{Haar}\bigr),
\label{open:bit-forcing}
\end{equation}
with the exchanged statements in the model sector. The solution exists for all $t\ge0$, and all bit and model means tend to zero. Rows, block rows, and the bit forcing are invariant under independent rotations of every belief and every model angle.
\begin{proof}
In \cref{open:full-field}, $2\log q_i=2\log\mu_i+2\log f_i$, so $h_i^q=2\log f_i(\theta\mid z)+\phi_i(z)$, where $\phi_i$ collects all terms depending on the bit alone. Integrating $\partial_tq_i=-q_i(h_i^q-\E_{q_i}h_i^q)$ over the angle gives $\partial_t\log\mu_i(z)=-(\E[h_i^q\mid z]-\E_{q_i}h_i^q)$. Subtracting, $\partial_t\log f_i=-2(\log f_i-\E[\log f_i\mid z])$, a separate Fisher-Rao relaxation toward Haar measure for each $z$, whose solution is \cref{open:conditional-flow}. Since $\E[h_i^q\mid z]=c_i(z,t)+\phi_i(z)$, the bit part reproduces the binary field and adds $-\frac12\mu_i(1)\mu_i(-1)[c_i(1,t)-c_i(-1,t)]$ to $\dot x_i$, with $\mu_i(1)\mu_i(-1)=1-x_i^2$. In logit coordinates the forcing is bounded by $\sup_z c_i(z,0)$, because the relative entropy decreases along \cref{open:conditional-flow}; the extremum comparison of \cref{open:interior} with this bounded forcing keeps the logits bounded on every finite interval, and since $c_i(z,t)\to0$ it gives convergence to zero. The explicit solution satisfies \cref{open:full-field} and therefore coincides with its unique local solution. Rows depend only on bit marginals, and $c_i$ is rotation invariant.
\end{proof}

The conditional relaxation \cref{open:conditional-flow} is the explicit solution of the Fisher-Rao gradient flow of a relative entropy \citep{ChenEtAl2023}, and the reduction is the Haar-reference analogue of \cref{smooth:split-dynamics}, with relaxation rate two. The receiver term $H_i^q=\KL(q_i\|\mu_i\otimes\mathrm{Haar})$ is the $\mu_i$-average of $c_i(z,t)/2$. For product laws $c_i(1,t)=c_i(-1,t)$ and the forcing vanishes. In particular, for arbitrary positive angular profiles, the product reference family
\begin{equation}
q_i(z,\theta)=\mu_i(z)f_i(\theta),\qquad
s_i(w,\psi)=\eta_i(w)g_i(\psi)
\label{open:product-reference}
\end{equation}
is invariant, and its angular factors obey
\begin{equation}
\partial_tf_i=-2f_i(\log f_i-\E_{f_i}\log f_i),
\qquad
f_i(\theta,t)=
\frac{f_i(\theta,0)^{e^{-2t}}}
{\int f_i(\psi,0)^{e^{-2t}}d\psi/(2\pi)}.
\label{open:angular-flow}
\end{equation}
A von Mises factor $p_\kappa(\theta-a)=e^{\kappa\cos(\theta-a)}/I_0(\kappa)$ keeps its phase and has concentration $\kappa e^{-2t}$.

\propositionheading{Open-law transfer}{open:open-laws}
The routing behavior of \cref{open:routing} persists on a nonempty relabeling-invariant set of independently variable positive belief and model densities on $\{-1,1\}\times S^1$, including correlated bit-angle laws. The set is open for uniform convergence of log-densities, and hence in the $C^r$ chart.
\begin{proof}
Let $M$ bound the oscillation of every log-conditional $\log f_i(\cdot\mid z)$. For $\ell=\log f$ and $\rho=e^{-2t}$ write $\mathcal R_\rho(\ell)$ for the relative entropy to Haar measure of the law proportional to $e^{\rho\ell}$. Its derivative in a direction $\varsigma$ is $\rho^2\operatorname{Cov}(\ell,\varsigma)$ under that law, which is at most $\rho^2M\|\varsigma\|_\infty$ in absolute value along the segment between the two log-conditionals. Hence
\begin{equation}
|c_i(1,t)-c_i(-1,t)|\le2e^{-4t}M\bigl\|\log f_i(\cdot\mid1,0)-\log f_i(\cdot\mid-1,0)\bigr\|_\infty.
\end{equation}
In the scaled coordinates of \cref{open:routing}, the forcing \cref{open:bit-forcing} is therefore a time-dependent field perturbation of size at most $e^{-4t}M\epsilon/h$, where $\epsilon$ is the largest of these conditional differences over agents and sectors. Take the binary open set of \cref{open:routing}, shrunk so that its windows survive a field error $e_0>0$ by \cref{open:persistence}. Every law whose bit and model means lie in it and whose conditionals satisfy $M\epsilon<he_0$ has the routing behavior. This condition is open for uniform convergence of log-densities and contains every product law \cref{open:product-reference}; rows and block rows are the binary ones by \cref{open:reduction}. Admissible bit-angle correlations are thus of order $h$.
\end{proof}

\corollaryheading{Transient relations and a forced final merge}{open:final-merge}
For every initial law of \cref{open:open-laws}, every native row tends to $1/N$. Every first-stage contrast tends to $1/6-4/6<0$ and the core-parent contrast against the third pair tends to $0<\delta$, so no block relation used in \cref{open:routing} is forward invariant. After the core-parent join the partition has two blocks, the exterior sum in \cref{open:contrast} is empty, and the remaining block rows tend to $2/6$ and $4/6$. Under \cref{open:selector} every such trajectory therefore reaches the one-block partition at a finite time after $L$, with an empty exterior. Under the retention below, this last join discards one further relative orientation, since $\mathsf H_{I\cup J}\times\mathsf H_K\subsetneq\mathsf H_{\{1,\ldots,6\}}$.
\begin{proof}
All means tend to zero by \cref{open:reduction}, so every divergence in \cref{open:row} tends to zero. With two blocks the contrast condition reduces to $B^a>\varepsilon$ in both directions and sectors. The cross block rows are below $10^{-3}$ at the core-parent join and tend to limits above $\varepsilon$, so there is a finite time after which every strict inequality holds, and the event occurs no later than that time.
\end{proof}

The persistence in this chapter is therefore finite-window by construction. The integrations of \cref{open:null-control} show the same final merge along the null limit flow.

\section{Orientation compression and inherited dynamics}

For an angle $a$, let $(\tau_af)(z,\theta)=f(z,\theta-a)$. An active rotation of agent $i$ applies $\tau_{a_i}$ to its belief and model laws in fixed sample coordinates. It changes the laws, unlike the passive changes of presentation of \cref{ch:interactions}. The energy, metric, rows, and \cref{open:full-field} are invariant or equivariant under independent active rotations of every belief and every model angle, a group $(S^1)^{2N}$. Below we use the diagonal subgroup $(S^1)^N$ that rotates both laws of one agent by the same angle. Whether an active rotation is detectable is a statement about a declared observable algebra (\cref{open:observables}): no agent, channel, row, or evaluator of the population resolves it, and only the external response protocol of the next section can.

For a block $I$ discard the subgroup
\begin{equation}
\mathsf H_I=\{(a_i)_{i\in I}:\sum_{i\in I}a_i=0\pmod{2\pi}\}.
\label{open:physical-group}
\end{equation}
Its dimension is $|I|-1$. A quotient alone does not give the required contextual statistical realization, so retain a whole local contextual section with fixed anchor $c_*$. Suppose $u_i(c_*)=\E_{q_i(c_*)}e^{i\theta}\ne0$, choose compatible local phase lifts, and define
\begin{equation}
\vartheta_i=\arg u_i(c_*),\qquad
\phi_I=\frac1{|I|}\sum_{i\in I}\vartheta_i,\qquad
a_i=\phi_I-\vartheta_i.
\label{open:anchor}
\end{equation}
This normalization is a first-Fourier-mode slice \citep{BudanurEtAl2015} in the method of slices for equivariant evolution equations \citep{RowleyEtAl2003,BeynThummler2004}, a cross-section in the sense of moving frames \citep{FelsOlver1998}; it is singular where an anchor moment vanishes. Apply $a_i$ to the entire section, independently of evaluation context:
\begin{equation}
\widetilde q_i(c)=\tau_{a_i}q_i(c),\qquad
\widetilde s_i(c)=\tau_{a_i}s_i(c),\qquad
q_I=\bigotimes_{i\in I}\widetilde q_i,\quad
s_I=\bigotimes_{i\in I}\widetilde s_i.
\label{open:parent-law}
\end{equation}
These are actual normalized parent laws. Their factors are recovered by marginalization. The relative phase profiles $\arg u_i(c)-\arg u_i(c_*)$ remain retained; only constant relative offsets are discarded. On this chart the fibers are precisely the constant $\mathsf H_I$ orbits. Choosing the offsets separately at each context would instead remove contextual phase variation itself and change the contextual geometry; the anchored choice avoids this. The model offsets are rotated by the belief-anchored angles. Independent model offsets are also symmetries and could be discarded as well; we do not do so.

The parent retains every child law up to one constant rotation per child with zero sum, together with its primitive ports. Its state size grows with the number of children, and all public bit marginals are retained. In the sense of the closing discussion of \cref{ch:smooth-closure}, it does not combine distinct public child laws into a smaller collective state.

Parent references and evaluators are products of their normalized primitive data. For each primitive port $i\in I,j\in J$, project a source-parent sample onto $j$, apply \cref{open:channel}, and compare with factor $i$. This is a normalized inherited channel. A full-parent version appends independent recipient factors on the other coordinates:
\begin{equation}
\mathcal L_{IJ}^{ij}(\boldsymbol k_J,d\boldsymbol k'_I)
=\mathcal K(k_j,dk'_i)
\bigotimes_{\ell\in I\setminus\{i\}}\widetilde q_\ell(dk'_\ell).
\end{equation}
Consequently $\KL(q_I\|q_J\mathcal L_{IJ}^{ij})=D^q_{ij}$ by cancellation of the matching normalized factors. If this representation is differentiated, its recipient-fill channel derivatives must be retained. It represents the original primitive sums and rows, not one newly fitted parent KL cost.

At a join of $I,J$, rotate their factors further by $\phi_{I\cup J}-\phi_I$ and $\phi_{I\cup J}-\phi_J$, where
\begin{equation}
\phi_{I\cup J}=
\frac{|I|\phi_I+|J|\phi_J}{|I|+|J|}.
\label{open:phase-recursion}
\end{equation}
Each primitive total rotation becomes $\phi_{I\cup J}-\vartheta_i$, so direct and sequential retention agree. Primitive counts are inherited. Since
$\mathsf H_I\times\mathsf H_J\subsetneq\mathsf H_{I\cup J}$,
this join discards one further continuous relative block orientation. Native row trajectories and all selected event times are unchanged by these rotations, so recording the times does not recover the discarded offsets.

\remarkheading{Global representatives}{open:global-anchor}
For $n=|I|\ge2$ no map $(S^1)^n\to S^1$ is both invariant under relabeling and equivariant under a common rotation. If $v$ is the rotation by $2\pi/n$ and $x=(e,v,\ldots,v^{n-1})$, then $vx$ is a cyclic relabeling of $x$, so such a map would give $\phi(x)=\phi(vx)=v\phi(x)$. The equal-phase anchor \cref{open:anchor} therefore exists only on local phase charts, with or without continuity. Dropping relabeling symmetry removes the obstruction: rotating every agent $j\ge2$ of $I$ to anchor phase zero and the first agent by the sum of their phases is a continuous $\mathsf H_I$-canonical representative on the whole domain where the anchor moments are nonzero. Rows, events, response laws, and collective costs are $\mathsf H_I$ invariant, so they do not depend on which representative is used. The common rotation meets $\mathsf H_I$ in $\mathbb Z_n$; on parent orbits the common group acts through $S^1/\mathbb Z_n$, and the parent sees the holonomy of the common connection only modulo $\mathbb Z_n$.

\propositionheading{Original-time parent flow}{open:parent-flow}
On the regular local chart the retained network of sections has an autonomous flow intertwining the original dynamics at the same time.
\begin{proof}
Equivariance and uniqueness imply that two sections related by a constant $\mathsf H_I$ rotation remain so for their common existence interval. This defines the quotient flow, the projectability condition underlying exact lumpability \citep{HorstmeyerAtay2016}. An explicit representative equation is also available. Let
\begin{equation}
\omega_i=\operatorname{Im}
\frac{\int e^{i\theta}V_i^q(c_*)d\nu}{u_i(c_*)},
\qquad
\overline\omega_I=|I|^{-1}\sum_{i\in I}\omega_i.
\end{equation}
This ratio is invariant under rotation and is computable from retained factors. Since $\dot a_i=\overline\omega_I-\omega_i$ and $\partial_a\tau_af=-\partial_\theta\tau_af$,
\begin{equation}
\partial_t\widetilde q_i
=V_i^q(\widetilde x)
-(\overline\omega_I-\omega_i)\partial_\theta\widetilde q_i,
\label{open:moving-flow}
\end{equation}
with the same coefficient for $\widetilde s_i$. This is the in-slice equation of the freezing method \citep{BeynThummler2004}. The product rule gives the parent-law derivative. Every term is retained. Successive correction coefficients telescope by \cref{open:phase-recursion}. For finite $C^r$ regularity the displayed derivative is in $C^{r-1}$; existence and uniqueness follow from the fine flow and equivariance, without assuming rotations are differentiable as maps $C^r\to C^r$.
\end{proof}

The Fisher metric of this canonical product slice need not be the kinetic quotient metric. For two location families with inertias $I_1,I_2>0$ and retained total phase $F$, minimizing lift norm subject to $\dot\vartheta_1+\dot\vartheta_2=\dot F$ gives
\begin{equation}
ds_{\rm quotient}^2=\frac{dF^2}{I_1^{-1}+I_2^{-1}},
\qquad
ds_{\rm slice}^2=\frac{I_1+I_2}{4}dF^2.
\label{open:two-metrics}
\end{equation}
They differ unless the inertias agree. The first is the quotient metric of a Riemannian submersion \citep{ONeill1966}; the second is the metric induced on the section. The exact parent update is \cref{open:moving-flow}, not an uncorrected unit-Fisher gradient of a new single parent KL. Its gradient structure is that of the quotient metric.

\propositionheading{The parent flow and the quotient metric}{open:quotient-gradient}
Let $V$ be the fine field of \cref{open:full-field} and let $\overline g$ be the quotient metric on $\mathsf H_I$ orbits, $\overline g(\pi_*v,\pi_*v)=\min_w g(v+w,v+w)$ over orbit directions $w$, where $g$ is the unit Fisher metric. At every law of the regular chart, $V$ is Fisher-orthogonal to the $\mathsf H_I$ orbit directions, the moving term in \cref{open:moving-flow} is the orbit component of the slice velocity, and the descended potential $\overline{\mathcal E}$ satisfies
\begin{equation}
\frac{d}{dt}\overline{\mathcal E}=-\overline g(\pi_*V,\pi_*V).
\label{open:quotient-dissipation}
\end{equation}
On the finite-dimensional invariant family of products of binary marginals and von Mises angular factors with positive concentrations, the $\mathsf H_I$ action is free, proper, and Fisher-isometric, and the parent flow is the gradient flow of $\overline{\mathcal E}$ for $\overline g$ at the original time.
\begin{proof}
The potential is $C^1$ for the supremum norm of positive densities, and for a density of class $C^1$ in the angle the orbit curve $a\mapsto\tau_aq$ is $C^1$ in that norm with tangent $-\partial_\theta q$. Invariance gives $d\mathcal E(J\xi)=0$ for every orbit direction $J\xi$, hence $g(V,J\xi)=-d\mathcal E(J\xi)=0$. The slice velocity in \cref{open:moving-flow} is $V$ plus an orbit direction. Since $V$ is orthogonal to orbits, $\overline g(\pi_*V,\pi_*V)=g(V,V)$, and $\frac{d}{dt}\overline{\mathcal E}=d\mathcal E(V)=-g(V,V)$. The family of the last statement is invariant by \cref{open:angular-flow}; the full field is tangent to it, so it equals the Fisher gradient of the restricted potential for the induced metric. A rotation fixing a von Mises factor with positive concentration is trivial, the torus is compact, and rotations preserve the Fisher metric and the potential. The orbit map is then a Riemannian submersion for $\overline g$, and \cref{inh:gradient-inheritance} with $d=1$ and $\chi=\id$ gives same-time inheritance of the gradient flow.
\end{proof}

On the finite-$r$ chart a rotation is not differentiable as a map $C^r\to C^r$, so no quotient Banach manifold is asserted there; the pointwise dissipation identity \cref{open:quotient-dissipation} is what holds on the whole chart. On the von Mises family the phases are constant along the flow, so the moving term vanishes there and the last statement does not exercise it.

\section{Responses and retained histories}

Declare one recursive response rule before any selection. A query to current block $I$ draws independent constituent samples $(Z_i,\Theta_i)\sim q_i$ and returns
\begin{equation}
\left((Z_i)_{i\in I},\ \sum_{i\in I}\Theta_i\pmod{2\pi}\right).
\label{open:response}
\end{equation}
The model sector has its corresponding independent joint response. Native source draws, transmitted bit samples, and joint measurable functions of whole-block outputs are also admitted. A policy may depend on past records, but cannot split a compressed block to reveal an individual absolute angle. Queries are observational and do not reset the law state. The querier is not an agent of the population and does not act on its dynamics; it reads angles in a fixed angular frame that no agent holds, since every channel twirls the angle.

\definitionheading{Declared observable algebra}{open:observables}
For a current partition $\mathcal P$, let $\mathcal A_{\rm pop}$ be the algebra of bounded measurable functions of the configuration generated by the quantities the population reads: rows, channel output laws, evaluators, and costs. The vector field is equivariant, so the population's evolution of these quantities does not depend on the orbit representative. Let $\mathcal A_{\mathcal P}$ be generated by $\mathcal A_{\rm pop}$ and by the probabilities of finite adaptive response histories under \cref{open:response} for $\mathcal P$. A distinction between two configurations is resolved by an algebra if some element takes different values on them.

Every element of $\mathcal A_{\rm pop}$ is invariant under $(S^1)^{2N}$, by \cref{open:reduction} and the invariance of the potential. Thus no active rotation is resolved within the population, and the words gauge and physical have no content here without the response protocol. The kernel in \cref{open:response} is a normalized pushforward of a product law. Since $\sum_ia_i=0$, it is exactly the same pushforward of the parent law in \cref{open:parent-law}. Its first angular moment is $\prod_iu_i$, generally nonzero, so an aggregate angle is resolved by $\mathcal A_{\mathcal P}$. Joining concatenates bit vectors and adds angle sums. Thus the same experiment composes at each level.

Equal retained states have equal waiting-time evolution and equal next-record joint kernels. Conditioning successively on the realized history proves equality of every finite adaptive joint-history law. Prior records are carried forward. A finite collection of sampled angular records does not deterministically identify the erased offsets: all the densities are positive, leaving continuously many compatible configurations for the same finite history. An exact record of constituent laws would instead invalidate the loss assertion. An individual absolute-angle query also fails descent immediately because its first moment changes under \cref{open:physical-group}. These experiments are outside the declared interface after the handoff.

\remarkheading{How much is lost}{open:loss-information}
The loss is non-identifiability, not absence of information. Before a merge every agent is its own block, so \cref{open:response} returns samples of its angular marginal. For an angular marginal $f$, the Fisher information about its phase carried by one sample is the rotation inertia $\E_f[(\partial_\theta\log f)^2]$, which is $\mathcal I(\kappa)$ below for a von Mises factor. On the product reference family the phases are constant, so $n_i$ pre-merge records of agent $i$ carry information $n_i$ times its inertia about the offsets that the merge discards, and their circular mean estimates those offsets with error of order $(n_i\E_f[(\partial_\theta\log f)^2])^{-1/2}$. This information is finite for every finite record set and grows linearly with the number of records. A uniform loss statement needs a protocol that bounds the number of pre-merge angle queries or discards records at merges; the present protocol does neither.

\section{Collective geometry on actual parent laws}

Take the common collective group to be $S^1$, acting on every belief and model angle simultaneously. This tie is a declared representation. The constitutive data are also invariant under independent belief and model rotations; for the untied group $S^1_q\times S^1_s$ the minimizing connection below is $A_q=-d\chi$, $A_s=0$, which is flat, and the minimized tensor has rank one. The curvature computed below belongs to the tied representation. No row, channel, record, or vector field of this construction reads the connection; $A^F$ is a diagnostic of the evolved sections.

The common rotation commutes with the independent active rotations. Use the original sum Fisher metric and convention $D^A p=dp+J_pA$. At fixed time the constants $a_i$ in \cref{open:anchor} do not vary with evaluation context. Hence $d\widetilde q_i=\tau_{a_i}dq_i$ and similarly for $s_i$. A product score is a sum of centered independent factor scores, so all cross terms vanish. Rotation preserves the remaining integrals. For every contextual direction $X$ and scalar gauge velocity $\xi$ this proves
\begin{equation}
\overline g(dRpX+\overline J\xi,dRpX+\overline J\xi)
=g(dpX+J\xi,dpX+J\xi).
\label{open:cost-equality}
\end{equation}
Equality holds before minimization. The orbit inertia, contextual momentum, minimizing connection class, and minimized tensor are therefore identical. On a regular chart their curvature agrees. This is a calculation from normalized laws, rather than assigning the fine tensor to retained metadata. It is consistent with fixed-statistic Fisher sufficiency results \citep{AyJostLeSchwachhofer2015}, but the present retention is section valued.

Under a passive common rotation $\chi_0(c)$, each anchor phase shifts by $\chi_0(c_*)$, leaving every $a_i$ unchanged. Parent sections transform by the same $\chi_0(c)$, and $A'=A-d\chi_0$ gives the usual covariant derivative. The local Lie-algebra stabilizer is preserved. Globally the total phase has $|I|$ roots; a common rotation $2\pi/|I|$ lies in the discarded group (\cref{open:global-anchor}). The construction stays on compatible local charts.

A nonflat family accompanies the routing trajectory. Define
\begin{equation}
p_\kappa(\theta)=\frac{e^{\kappa\cos\theta}}{I_0(\kappa)},
\qquad
\mathcal I(\kappa)=\kappa\frac{I_1(\kappa)}{I_0(\kappa)}.
\end{equation}
Here $I_0$ normalizes Haar measure and $I_1/I_0=\E_{p_\kappa}\cos\theta$ \citep{MardiaJupp2000}.
Integrating the derivative of $\sin\theta p_\kappa(\theta)$ around the circle gives rotation inertia
$\E_{p_\kappa}[\kappa^2\sin^2\theta]=\mathcal I(\kappa)$.
For $\kappa>0$ this mean cosine is positive, and its derivative is
$\operatorname{Var}_{p_\kappa}(\cos\theta)>0$, so $\mathcal I,\mathcal I'>0$.

Use contexts $(\kappa,\chi)$, belief factors $p_\kappa(\theta-\chi)$, and model factors $p_\lambda(\psi)$ with fixed $\lambda>0$. Binary coordinates are initially context independent. Put $\rho_t=e^{-2t}$,
$a_t=\mathcal I(\rho_t\kappa)$ and $b_t=\mathcal I(\rho_t\lambda)$. Parity makes concentration and rotation scores orthogonal. For a block of primitive count $n$ the common-gauge cost in its angular direction is
$n[a_t(d\chi+\xi)^2+b_t\xi^2]$. Thus
\begin{equation}
A_t^F=-\frac{a_t}{a_t+b_t}d\chi,\qquad
F_t^F=-\frac{b_t\rho_t\mathcal I'(\rho_t\kappa)}
{(a_t+b_t)^2}d\kappa\wedge d\chi,
\label{open:curvature}
\end{equation}
and
\begin{equation}
h_{\chi\chi}=n\frac{a_tb_t}{a_t+b_t}>0,\qquad
h_{\kappa\kappa}
=n\rho_t^2\operatorname{Var}_{p_{\rho_t\kappa}}(\cos\theta)>0,
\qquad h_{\kappa\chi}=0.
\label{open:rank}
\end{equation}
Curvature is nonzero at every finite time. As $t\to\infty$, $\mathcal I(v)\sim v^2/2$ and $\mathcal I'(v)\sim v$ give
\begin{equation}
F_t^F\longrightarrow-\frac{2\lambda^2\kappa}{(\kappa^2+\lambda^2)^2}d\kappa\wedge d\chi,
\qquad
e^{4t}h_{\chi\chi}\longrightarrow\frac{n\kappa^2\lambda^2}{2(\kappa^2+\lambda^2)},
\qquad
e^{4t}h_{\kappa\kappa}\longrightarrow\frac n2.
\label{open:geometry-limits}
\end{equation}
On the reference family the curvature coefficient is therefore bounded away from zero uniformly in $t\ge0$ on compact context patches with $\kappa>0$, while the tensor decays like $e^{-4t}$; the infinite-time degeneration is a loss of rank, not of curvature. The curvature is already present at $t=0$ and does not depend on the selected partition. The retention preserves it because \cref{open:cost-equality} holds for every partition at every time (\cref{open:all-partitions}). This geometry is supplied by the declared contextual family; it is not produced by the selection.

To extend this geometry through every nearby law value without encoding its erased orientation, fix $c_*=(\kappa_*,0)$. For arbitrary nearby initial $x=(q_i,s_i)_i$, let $\vartheta_i(x)=\arg\E_{q_i}e^{i\theta}$ and assign
\begin{equation}
q_{i,x}(\kappa,\chi)
=\tau_\chi\left[
\frac{e^{(\kappa-\kappa_*)\cos(\theta-\vartheta_i(x))}q_i(z,\theta)}
{\int e^{(\kappa-\kappa_*)\cos(\theta-\vartheta_i(x))}q_i(z,\theta)d\nu}
\right],
\qquad s_{i,x}(\kappa,\chi)=s_i.
\label{open:equivariant-lift}
\end{equation}
This normalized section passes through $x$ at the anchor and reproduces the preceding family at the product reference. An active rotation shifts both the density and its centering phase, so the entire section rotates by that same constant. Thus its retained profiles do not secretly store the discarded offsets.

On a compact context patch with $\kappa>0$ and a finite time horizon, the reference densities, moments, inertia, tensor eigenvalues, and absolute curvature coefficient have positive lower bounds. Smooth dependence of the fine field and its contextual variational equations preserves these bounds on a sufficiently small full neighborhood of independent initial laws; this step uses the $C^r$ chart, since the rotation scores are angular derivatives. The lift chooses the contextual jets on which geometry is asserted; it imposes no product or common-parameter restriction on their anchor values.

Select membership at the fixed anchor and hand off the whole section on that common time slice. This convention matters. If each context instead chooses its own birth time $\tau(c)$, the derivative is
\begin{equation}
d_c[p_{\tau(c)}(c)]X
=(d_cp_tX)_{t=\tau(c)}+V(p_{\tau(c)}(c))X\tau.
\label{open:event-derivative}
\end{equation}
That distinct comparison includes the additional score and is not implied by \cref{open:cost-equality}.

\section{The combined existence statement}

The inheritance statements above never use the selection. This is the content of the following proposition, which makes the division of labor in the main theorem explicit.

\propositionheading{Symmetry reduction for every partition}{open:all-partitions}
Let $\mathcal P$ be any partition of the agents, let $t_0\ge0$, and suppose every anchor moment is nonzero at $t_0$. Put $\mathsf H_{\mathcal P}=\prod_{I\in\mathcal P}\mathsf H_I$. Then every row, block row, selector input, and event time is invariant under $(S^1)^{2N}$, hence under $\mathsf H_{\mathcal P}$; the anchored retention of $\mathcal P$ at $t_0$ has the original-time flow of \cref{open:parent-flow}; every finite adaptive history of \cref{open:response} for the current partition $\mathcal P$ descends; and the collective minimizing costs satisfy \cref{open:cost-equality}.
\begin{proof}
The first statement is \cref{open:reduction}. The proofs of \cref{open:parent-flow}, of response descent, and of \cref{open:cost-equality} use only equivariance under $\mathsf H_I$, the anchor, and the product form of the parent laws, and none uses an event time or a selected block.
\end{proof}

The selected partition is therefore not distinguished by inheritance. Native routing contributes a law-determined and relabeling-natural choice of partition and of times. The response rule contributes the coarsening that makes the discarded directions unresolvable. Before an event the current partition is finer: a singleton query returns $\Theta_i$, whose first moment changes under a rotation in $\mathsf H_{\{i,j\}}$, and a block query returns a sum whose first moment changes under a rotation in $\mathsf H_{I\cup J}$ that is not in $\mathsf H_I\times\mathsf H_J$. After the event these responses are no longer admitted. The distinction discarded at a handoff is exactly the one the declared interface stops resolving at that handoff; the handoff loss is created by the response interface, not by the selection.

\theoremheading{Open-law two-stage selection with exact inheritance}{open:formation}
Consider the reciprocal law potential \cref{def:dyn-potential} with the constitutive data of this chapter: six agents, bit-circle sample spaces, the selective channel \cref{open:channel} on every edge including self, uniform references and evaluators, uniform all-source priors, unit Fisher mobilities, and one fixed temperature. There exist $T>0$, a finite horizon, and a nonempty relabeling-invariant set of independently variable positive belief and model laws, open in the $C^r$ chart, on which the following hold over that horizon and on compatible local phase charts.
\begin{enumerate}
\item[(a)] Selection. \Cref{open:selector} selects three pairs, which coexist over a positive window, and then joins two of their parents while the third pair is a nonempty exterior with strictly smaller routing masses. Rows and events depend only on the bit marginals.
\item[(b)] Interface coarsening. At each event the response interface \cref{open:response} passes to the coarser partition. A continuous distinction between laws, along the newly discarded subgroup, is resolved before the event by admitted queries and after it by no admitted query; no element of $\mathcal A_{\rm pop}$ resolves it at any time. Records acquired before the event carry finite information about it (\cref{open:loss-information}).
\item[(c)] Exact inheritance. For the selected partitions, the anchored section-valued parent laws, inherited port channels, original-time parent flow, and every finite adaptive history of \cref{open:response} descend exactly.
\item[(d)] Geometry. On the declared contextual family \cref{open:equivariant-lift}, the handoffs preserve the tied collective minimizing connection and a rank-two tensor with nonzero curvature.
\end{enumerate}
The temperature is not quantified (\cref{open:constants}). The theorem does not establish interaction necessity, and the block relations it selects are transient (\cref{open:null-control,open:final-merge}).
\begin{proof}
Use \cref{open:routing} to fix finite $A,R,h$ and its horizon, then \cref{open:open-laws} to obtain a full positive-law neighborhood. Choose the angular profiles with nonzero first moments and the equivariant sections of \cref{open:equivariant-lift}. Shrink this neighborhood in the $C^r$ chart to preserve the geometric bounds, phase chart, and all routing windows. Intersection remains nonempty because the same product trajectory satisfies them with strict compact margins.

The primitive laws are distinct and stay distinct at finite regular times by uniqueness on the invariant equality diagonals. The selector and inherited rows are exactly \cref{open:selector,open:row-recursion}, and (a) follows. The first joins discard their $\mathsf H_I$ directions; the later join discards an additional relative parent direction. Small active rotations stay within the open neighborhood, and equivariance of the assigned sections preserves these fibers after all retained profile data are included. Rows and event times are invariant along them. Statement (b) is the paragraph after \cref{open:all-partitions}, together with the invariance of $\mathcal A_{\rm pop}$ and positivity of the anchor moments. Parent normalization, port channels, recursive composition, original-time flow, and history descent in (c) follow from \cref{open:parent-law,open:phase-recursion,open:parent-flow} and \cref{open:all-partitions}. Equation \eqref{open:cost-equality} preserves the actual minimizing costs, and \cref{open:curvature,open:rank} establish their nontrivial geometry throughout the same finite interval, proving (d).
\end{proof}

Against \cref{inh:recursive-formation} the theorem gives (F1) with the $C^r$ topology, (F2) with finite-window persistence, (F3) with the response class declared before selection, and (F4), since each retention is non-injective and a pre-handoff query distinguishes its fibers. It does not establish (F5) (\cref{open:null-control}). The first part of (F6) holds: the parent pair consists of first-stage parents under inherited rows, without reset or retuning. Its counterfactual part is not verified, since the second stage depends on the first only through the partner exclusion of \cref{open:selector} and occurs when the supplied core separation has contracted to order one. The construction is therefore a two-stage selection with exact inheritance in the sense of the standard, with that qualification, and not a recursive formation. The symmetry enabling compression is fixed in the selective architecture, and the parent is not a fixed-size ordinary KL agent. The result gives neither a preferred threshold nor generic selection for arbitrary channels or all initial laws.

\remarkheading{Finite depth by collapsed clusters}{open:depth}
In \cref{open:limit}, two clusters of $m$ coincident agents at $\pm a$ have inherited cross mass $B=v/(1+v)$ with $v=e^{-2a^2}$, and $\dot a=-a(1+4B)$, independently of $m$: each member's row has $m$ unit entries and $m$ entries $v$. The late comparison in the proof of \cref{open:routing} therefore applies to collapsed clusters of every size. An induction over geometric scales would plausibly give any finite number of selected stages on a finite horizon, at a temperature that decreases with the depth. The uniformity of the tube estimates in the depth and the effect of the self mass on the contrast at large depth have not been checked, and we do not claim the result.

\openproblemheading{Interaction-necessary selection in the bit-circle architecture}{open:necessity-open}
Prove a version of \cref{open:formation} on an open set where the null flow does not reproduce the membership history, so that (F5) holds; integrations of the limit field at moderate separations suggest that such sets exist. Separately, find an architecture with a forward-invariant set on which a block relation holds and a later join of first-stage parents occurs; a compression symmetry produced by the dynamics rather than fixed by the channel, for example through an angle-transmitting channel of small concentration that makes the compression approximate on the selected blocks with a computable defect; a loss statement in which the information about discarded coordinates given all admitted records stays bounded; and explicit or validated witness constants.

\section{Two tests of stronger inheritance claims}

The first test makes the identity-channel response obstruction quantitative. For two binary source means $x_1\ne x_2$ and fixed positive weights, an exterior receiver mean $r$ determines the block normalizer. With
$a_j=\operatorname{atanh}x_j$, $v_j=a_j/T$, and
$\psi(r)=\frac{1+r}{2}\log(1+r)+\frac{1-r}{2}\log(1-r)$, removing its receiver entropy factor gives
\begin{equation}
F(r)=e^{\psi(r)/T}Z_I(r)=\sum_{j=1}^2 b_je^{v_jr},
\qquad b_j>0.
\end{equation}
At an interior $r_0$, set $m_k=F^{(k)}(r_0)$ and $c_j=b_je^{v_jr_0}$. Then
\begin{equation}
H=m_0m_2-m_1^2=c_1c_2(v_1-v_2)^2>0.
\end{equation}
The two scaled logits $v_j=a_j/T$ are the roots of
\begin{equation}
z^2-\frac{m_0m_3-m_1m_2}{H}z
+\frac{m_1m_3-m_2^2}{H}=0.
\end{equation}
This is Prony's method for a two-term exponential sum \citep{Prony1795}. Thus equality of response functions on any open receiver interval identifies the unordered source laws. Small within-pair KL makes this inverse ill-conditioned, not exactly redundant. This is a finite reconstruction within the unrestricted Gibbs-interface result; it does not obstruct the selective angular interface above.

The second test concerns propagation of Fisher sufficiency. Let a fixed statistic retain $Z$ from positive densities
$p_t(c;z,y)=\mu_t(c;z)\varkappa_t(c;y\mid z)$.
Write $P_t$ for conditional expectation given $Z$, let $X$ be a contextual direction, and suppose the complete centered force is
$\partial_t\log p_t=-(h_t-\E h_t)$.
The hidden contextual score is $a_X=X\log\varkappa_t$. Conditional normalization gives
\begin{equation}
\partial_ta_X=-(I-P_t)Xh_t+
\operatorname{Cov}_{\varkappa_t}(h_t,a_X).
\end{equation}
If $a_X(0)=0$, the Fisher loss therefore starts as
\begin{equation}
g_t(dp_tX,dp_tX)-\overline g_t(dRp_tX,dRp_tX)
=t^2\E_0[((I-P_0)Xh_0)^2]+o(t^2).
\label{open:generated-loss}
\end{equation}
This follows by differentiating the conditional score and expanding its squared norm; on finite positive alphabets all operations are finite sums.

For two identity-coupled laws $q_i(z,y)=\mu_i(z)\varkappa_i(y\mid z)$ on a finite alphabet, written as probabilities with respect to counting measure, with local references $r_i$, the complete forces are
$h_1=2\log q_1-\log r_1-\log q_2-q_2/q_1$ and its exchanged expression. Initially vary only $\mu_2$ with tangent $v$, keeping both conditionals fixed. Conditional centering of the incoming ratios gives the coefficient
\begin{equation}
\sum_z v_z^2\left[
\frac{\chi^2(\varkappa_2\|\varkappa_1)}{\mu_1(z)}
+\frac{\mu_1(z)^2}{\mu_2(z)^3}\chi^2(\varkappa_1\|\varkappa_2)
\right].
\end{equation}
Here the conditional divergences are evaluated at $z$. Taking
$\mu_1=\mu_2=(1/2,1/2)$, $v=(1,-1)$, and $\varkappa_1(y=1\mid z)=1/4$, $\varkappa_2(y=1\mid z)=3/4$ for both values of $z$ gives $32/3$. Initial score sufficiency is lost immediately. This example uses a trivial collective group; with a nontrivial action one must test the actual minimizing residual rather than infer optimized loss from an arbitrary raw score. The anchored retention avoids this failure through the exact equality of its full minimizing costs.
